\scanpage{84}
% archive foliation: 83
% lines 1--4
d'où
\[
\Gamma=\operatorname{Gal}(\overline Q/Q)\longrightarrow\widetilde{\unsure{\mathcal T|\Gamma}}_\nu
\]
L'hom correspondant
\[
\Gamma\longrightarrow\unsure{\mathfrak S_{\nu+1}|\mathcal S_{\nu+1}}
\]
% lines 5--6
se fait via
\[
\Gamma\unsure{\text{---}|\longrightarrow}\mu_{\nu+1}(C^*)\simeq \mathbb Z/_{\nu+1}\xrightarrow{can}\unsure{\mathfrak S_{\nu+1}|\mathcal S_{\nu+1}}\ \ldots
\]
% lines 7--12
L'intérêt de prendre $I$ comme
dessus est d'avoir une opération
de $\inserted{\mathbb Z/_{2}\cdot}\mu_{n\,Q}$ \deleted{\unsure{tous}} sur la situation
$(X_0/I_0/\unsure{Q|k})$ donc de $\inserted{\mathbb Z/_{2}\cdot}\mu_n(\overline Q)\ (\simeq D_n)$
\margin{$n=\nu+1$}
sur $X/I/Q$ -- en fait, le
% lines 13--15
produit $\frac12$ direct $\bigl(\Gamma\unsure{\times|\cdot}\deleted{\unsure{\mathbb Z/_{2}}}\cdot\underbrace{\mu_n(\overline Q)}_{\unsure{\in|\subset}\overline Q^{*}}\bigr)\unsure{=}$
opère sur $X/I/\overline Q$ donc
\deleted{\unsure{s'envoie} dans}
% lines 16--18
\[
(\Gamma\times \mathbb Z/_{2})\cdot\unsure{\mu_{n+1}|\mu_n}(\overline Q)\unsure{\Longrightarrow|\longrightarrow}\widetilde{\unsure{\mathcal T|\Gamma}}_\nu
\]
où l'hom $\mu_n(\overline Q)\longrightarrow\widetilde{\unsure{\mathcal T|\Gamma}}_\nu$
